\chapter{Data Modeling}
\label{chap:data_modeling}

This chapter details the data persistence layer of \textbf{BookSphere}. Given the polyglot nature of the architecture, the data modeling strategy is split between a \textbf{Document-Oriented} approach (\textit{MongoDB}) for content management and analytics, and a \textbf{Graph-Oriented} approach (\textit{Neo4j}) for social interactions and recommendation algorithms.

\section{Data Source and Volume}
To ensure the system was tested against realistic data volumes and complex interaction patterns, we implemented a hybrid data generation strategy using custom automation scripts. The core catalog and foundational rating data are sourced from the \textbf{Amazon Books Reviews Dataset} (providing real-world book metadata and user ratings). 

This real data is a user base of approximately 10,000 active accounts. The generation process automatically populates the polyglot persistence layer:


\begin{itemize}
    \item \textbf{Amazon Books Reviews Dataset ($\sim$5GB):} Provides the core catalog of books and real-world reviews spanning from 1996 to 2014.
    \item \textbf{Book-Crossing Dataset:} Used as a supplementary source for user ratings to increase data \textit{variety} and \textit{volume}.
\end{itemize}
\begin{table}[h]
    \centering
    \renewcommand{\arraystretch}{1.5}
    \begin{tabular}{|p{0.45\textwidth}|p{0.45\textwidth}|}
        \hline
        \textbf{MongoDB (Document Store)} & \textbf{Neo4j (Graph Store - CSVs)} \\
        \hline
        \begin{itemize}
            \item 37.1MB for \texttt{reviews}
            \item 27.7MB for \texttt{books}
            \item 12.4MB for \texttt{users}
            \item 3.5MB for \texttt{authors}
        \end{itemize} 
        & 
        \begin{itemize}
            \item 4.32MB for \texttt{Nodes (User, Book, Review...)}
            \item 24.9MB for \texttt{Relationships (LIKES, FOLLOWS...)}
        \end{itemize} \\
        \hline
    \end{tabular}
    \caption{Database Volumes by Collection.}
    \label{tab:db_volumes}
\end{table}
\section{Document Store: MongoDB}
MongoDB acts as the primary "Source of Truth" for the system. It stores the complete information about books, users, and reviews. The schema design prioritizes \textbf{Read Performance} through specific NoSQL patterns.

\subsection{Collections Overview}
The database \texttt{booksphere\_db} contains the following main collections:

\begin{enumerate}
    \item \textbf{books:} The central catalog entity.
    \item \textbf{users:} Stores authentication data and the user's library.
    \item \textbf{reviews:} Stores the full text of reviews. Linked to books and users via references.
    \item \textbf{authors:} Stores aggregated statistics about authors (e.g., total ratings).
\end{enumerate}

\subsection{Schema Design (JSON)}

\subsubsection{Books Collection}
We implemented the \textbf{Bucket Pattern} (\texttt{stats\_per\_year}) to make historical ranking queries instant, avoiding the need to scan millions of review documents. The \textbf{Snapshot Pattern} (\texttt{recent\_reviews\_snapshot}) allows the frontend to render the book details page with the latest reviews in a single query.
\\
\begin{lstlisting}[language=json, caption={Example Book Document}, label={lst:book_example}]
{
  "_id": "698db76b4a1f5c8fa49331a8",
  "title": "The Fellowship of the Ring",
  "publication_year": 2000,
  "description": "Begin your journey into Middle-earth... The inspiration for the upcoming...",
  "availability": "ACTIVE",
  "source": "amazon_master",
  
  "author": {
    "id": "698db76b4a1f5c8fa4933176",
    "name": "J. R. R. Tolkien"
  },
  
  "genres": [
    "Fiction"
  ],
  
  "external_ids": {
    "isbns": [
      "978-0547928210", 
      "..."
    ]
  },
  
  "recent_reviews_snapshot": [
    {
      "_id": "698db7874a1f5c8fa493ddcf",
      "user_id": "698db7854a1f5c8fa493cad6",
      "username": "User_278550",
      "rating": 98,
      "summary": "",
      "date": "2025-12-27T00:00:00.000+00:00"
    }
    // ... other four recent reviews
  ],
  
  "popular_reviews_snapshot": [
    {
      "_id": "698db7ba4a1f5c8fa4940742",
      "user_id": "698db7854a1f5c8fa493bd2c",
      "username": "User_AASZHOEVN5LZH",
      "rating": 82,
      "num_of_like": 29,
      "summary": "The legacy of \"The Hobbit\"",
      "date": "2014-02-26T01:00:00.000+00:00"
    }
    // ... other two popular reviews
  ],
  
  "stats_per_year": [
   {
      "year": 2011,
      "ratings_count": 5,
      "sum_rating": 90
    }
    // ... between years
    {
      "year": 2025,
      "ratings_count": 1,
      "sum_rating": 98
    }
  ],
  
  "review_ids": [
    "698db7874a1f5c8fa493ddcf",
    "698db7874a1f5c8fa493e262",
    "698db7874a1f5c8fa493e3b8",
    "698db7ba4a1f5c8fa494072c",
    "..." 
  ],
  
  "month_score": {
    "rating_count": 1,
    "sum_rating": 98,
    "Current_Month": "2025-12"
  }
}
\end{lstlisting}

\subsubsection{Users Collection}
The user document utilizes the \textbf{Subset Pattern} for the \texttt{bookshelf} to keep the document size manageable. Crucially, we maintain a \texttt{reviews\_year} array (embedded) which is optimized specifically for the\textbf{ Yearly Wrapped} aggregation, allowing the system to generate the report without joining the \texttt{reviews} collection.
\\

\begin{lstlisting}[language=json, caption={User Document Structure}, label={lst:user_json}]
{
  "_id": "ObjectId('65d1...')",
  "username": "User_12345",
  "email": "user@bx.com",
  "password_hashed": "...", 
  "status": "active",
  
  "bookshelf": [
    {
      "book_id": "ObjectId('65b3f...')",
      "status": "read",
      "added_at": "ISODate('2023-06-01...')",
      "title": "The Hobbit",
      "author": { "name": "J.R.R. Tolkien" },
      "genres": ["Fantasy"]
    }
  ],
  
  "reviews_year": [
    {
      "id": "ObjectId('99a1...')",
      "rating": 85,
      "book": "Dune"
    }
  ]
}
\end{lstlisting}


\subsubsection{Reviews Collection}
The reviews collection stores the full textual content of user evaluations. It serves as the primary source for the "read more" functionality and provides the data for both the \texttt{recent\_reviews\_snapshot} in the books collection and the \texttt{reviews\_year} array in the user documents.
\\
\begin{lstlisting}[language=json, caption={Example Review Document}, label={lst:review_example}]
{
  "_id": "698db7874a1f5c8fa493ddcf",
  "user_id": "698db7854a1f5c8fa493c8a5",
  "username": "Rusneus003",
  "rating": 83,
  "source": "amazon",
  "text": "This is, all-around, a pretty good lawn book. However, there are a few cases where the advice is a little outdated...",
  "summary": "Pretty good, although a few flaws",
  "created_at": "2022-04-26T02:00:00.000Z",
  "likes_count": 6,
  "is_banned": false,
  "book_snapshot": {
    "book_id": "698db76e4a1f5c8fa493a4ec",
    "title": "The Impatient Gardener's Lawn Book"
  },
  "author_snapshot": {
    "id": "698db76e4a1f5c8fa493a4ed",
    "name": "Jerry Baker"
  }
}
\end{lstlisting}


\subsubsection{Authors Collection}
The authors collection is designed to facilitate quick retrieval of author profiles and their complete bibliography. It utilizes the \textbf{Extended Reference Pattern} to list published books and the \textbf{Computed Pattern} to maintain real-time aggregate statistics.
\\
\begin{lstlisting}[language=json, caption={Example Author Document (Edgar Allan Poe)}, label={lst:author_example}]
{
  "_id": { "$oid": "698db76c4a1f5c8fa4936c8d" },
  "name": "Edgar Allan Poe",
  "status": "ACTIVE",
  "published_books": [
    {
      "_id": { "$oid": "698db76c4a1f5c8fa4936c8e" },
      "title": "Tales of Mystery and Imagination"
    },
    {
      "_id": { "$oid": "698db76c4a1f5c8fa49373bf" },
      "title": "Fall of the House of Usher and Other Stories"
    },
    {
      "_id": { "$oid": "698db76d4a1f5c8fa4937aef" },
      "title": "The Raven And Other Poems"
    },
    {
      "_id": { "$oid": "698db76e4a1f5c8fa49395df" },
      "title": "The Pit and the Pendulum"
    },
    {
      "_id": { "$oid": "698db76e4a1f5c8fa4939a9c" },
      "title": "The Narrative of Arthur Gordon Pym of Nantucket"
    },
    {
      "_id": { "$oid": "698db76e4a1f5c8fa493a884" },
      "title": "Edgar Allan Poe Reader"
    }
  ],
  "ratings_count": 14,
  "sum_ratings": 1234
}
\end{lstlisting}

\subsubsection{Authors Collection}
The authors collection is designed to facilitate quick retrieval of author profiles and their complete bibliography. It utilizes the \textbf{Extended Reference Pattern} to list published books and the \textbf{Computed Pattern} to maintain real-time aggregate statistics.

\begin{lstlisting}[language=json, caption={Example Author Document (Edgar Allan Poe)}, label={lst:author_example}]
{
  "_id": { "$oid": "698db76c4a1f5c8fa4936c8d" },
  "name": "Edgar Allan Poe",
  "status": "ACTIVE",
  "published_books": [
    {
      "_id": { "$oid": "698db76c4a1f5c8fa4936c8e" },
      "title": "Tales of Mystery and Imagination"
    },
    {
      "_id": { "$oid": "698db76c4a1f5c8fa49373bf" },
      "title": "Fall of the House of Usher and Other Stories"
    },
    {
      "_id": { "$oid": "698db76d4a1f5c8fa4937aef" },
      "title": "The Raven And Other Poems"
    },
    {
      "_id": { "$oid": "698db76e4a1f5c8fa49395df" },
      "title": "The Pit and the Pendulum"
    },
    {
      "_id": { "$oid": "698db76e4a1f5c8fa4939a9c" },
      "title": "The Narrative of Arthur Gordon Pym of Nantucket"
    },
    {
      "_id": { "$oid": "698db76e4a1f5c8fa493a884" },
      "title": "Edgar Allan Poe Reader"
    }
  ],
  "ratings_count": 14,
  "sum_ratings": 1234
}
\end{lstlisting}

\section{Graph Store: Neo4j}
Neo4j is utilized strictly for modeling relationships and enabling high-performance traversal queries. The graph schema is designed to be lightweight: nodes primarily contain IDs linking back to MongoDB, while relationships carry timestamp data to support temporal queries.

% IMAGE COMMENTED OUT TO PREVENT ERRORS
% \begin{figure}[h]
%    \centering
%    \includegraphics[width=0.8\textwidth]{neo4j_schema_diagram.png}
%    \caption{Neo4j Graph Schema: Relationships between User, Book, Review, and Author.}
%    \label{fig:neo4j_schema}
% \end{figure}

\subsection{Nodes and Labels}
\begin{itemize}
    \item \textbf{User:} Contains \texttt{mongoId}, \texttt{username}, \texttt{country}.
    \item \textbf{Book:} Contains \texttt{mongoId}, \texttt{title}, \texttt{year}.
    \item \textbf{Review:} Contains \texttt{mongoId}, \texttt{rating}, \texttt{createdAt}.
    \item \textbf{Author/Genre:} Lightweight nodes for structural categorization.
\end{itemize}

\subsection{Relationships}
Specific relationships are modeled to support the recommendation engine:

\begin{itemize}
    \item \texttt{(:User)-[:FOLLOWS \{since: ...\}]->(:User)}: the social graph backbone.
    \item \texttt{(:User)-[:LIKES \{timestamp: ...\}]->(:Book)}: a user like a book.
    \item \texttt{(:User)-[:POSTED \{timestamp: ...\}]->(:Review)}: a user posted a review.
    \item \texttt{(:User)-[:LIKES \{timestamp: ...\}]->(:Review)}: a user like a review of another user.
    \item \texttt{(:User)-[:LIKES \{timestamp: ...\}]->(:Author)}: a user like an author.
    \item \texttt{(:User)-[:LIKES \{timestamp: ...\}]->(:Genre)}: a user like a book's genre.
    \item \texttt{(:Review)-[:REFER\_TO]->(:Book)}:connect a review with the book reviewed.
    \item \texttt{(:Book)-[:BELONGS\_TO]->(:Genre)}: a book belongs to one ore more genres.
    \item \texttt{(:Author)-[:WROTE]->(:Book)}: a author is the writer of a certain book.
\end{itemize}

\section{Polyglot Persistence Strategy}
\label{sec:consistency}

Maintaining data integrity between a Document Store and a Graph Database requires a strictly defined synchronization strategy. BookSphere prioritizes Eventual Consistency for high-frequency social features (Likes, Views) to ensure system availability, while enforcing Strict Consistency for critical Admin and Moderation operations.

\subsection{Synchronization Matrix}
The following table details how operations are propagated across the two databases:


\subsection{Error Handling Strategy}
To ensure robustness, we implemented a specific error handling flow:
\begin{itemize}
    \item \textbf{MongoDB Failure:} If the write to the primary DB fails, the operation aborts immediately. No try-catch block is used, allowing the \texttt{@Transactional} annotation to handle the rollback automatically.
    \item \textbf{Neo4j Failure:} If the write to MongoDB succeeds but Neo4j fails, a specific \texttt{try-catch} block intercepts the error and performs a manual rollback on MongoDB (deleting the just-created entity) to prevent "phantom data" and maintain consistency.
\end{itemize}

